% !PRINTED 444
\weberpage{444}
% running head: 444  Eighteenth Chapter.  § 122.

Furthermore, if $p$ does not divide $2D$, then by § 85
\[
D^{\frac{p-1}{2}} \equiv (D, p) \pmod{p},
\]
and hence:
\weq{6}{\sqrt{D}^p \equiv (D, p)\sqrt{D} \pmod{p}.}
If therefore $\alpha$ denotes, as above, an integer of the field $\Omega$:
\[
\alpha = \frac{x + y\sqrt{\Delta}}{2},
\]
then, if one multiplies by $2Q$, writes $x$ for $Qx$, and denotes by $\alpha'$ the number conjugate to $\alpha$:
\[
2Q\alpha = x + y\sqrt{D}, \quad 2Q\alpha' = x - y\sqrt{D},
\]
and from this it follows by \gesperrt{Fermat}'s theorem and by (6):
\weq{7}{\begin{aligned} \alpha^p &\equiv \alpha, & (D,p) &= +1,\\ \alpha^p &\equiv \alpha', & (D,p) &= -1, \end{aligned} \quad (\mathrm{mod}\ p).}
From this it now follows, if $f(k, \sqrt{-m})$ is any number of the form (4), by application of the polynomial theorem:
\weq{8}{\begin{aligned} f(k,\sqrt{D})^p &\equiv f(k^p, \sqrt{D}), & (D,p) &= +1,\\ f(k,\sqrt{D})^p &\equiv f(k^p, -\sqrt{D}), & (D,p) &= -1, \end{aligned} \quad (\mathrm{mod}\ p),}
in which $k^p$ stands for $(k)^p$, and so has the meaning of an actual power.

3. If $\mathfrak{P}$ is a prime ideal in $\mathfrak{K}$ that divides $p$, then the norm of $\mathfrak{P}$ is a power of $p$:
\weq{9}{N(\mathfrak{P}) = p^f.}

The exponent $f$ is the \gesperrt{degree} of the prime ideal $\mathfrak{P}$. For every arbitrary integer $\xi$ of the field $\mathfrak{K}$ we have
\weq{10}{\xi^{N(\mathfrak{P})} \equiv \xi \pmod{\mathfrak{P}},}
and if conversely $p^f$ is the lowest power of $p$ for which the congruence
\weq{11}{\xi^{p^f} \equiv \xi \pmod{\mathfrak{P}}}
is satisfied by \gesperrt{every} number $\xi$, then $f$ is the degree of the ideal $\mathfrak{P}$ (Bd.\,II, § 167).

The degrees of the prime ideals contained in the prime $p$ are now to be determined.

Since $p$ does not divide the fundamental number of the field $\mathfrak{K}$, $p$ is not divisible by the square of a prime ideal.
